% Generated draft: requires mathematical and source review.
\documentclass[11pt]{article}
\usepackage[margin=25mm]{geometry}
\usepackage{fontspec}
\setmainfont{Latin Modern Roman}
\usepackage{amsmath,amssymb,mathtools,unicode-math}
\setmathfont{Latin Modern Math}
\usepackage{enumitem,needspace,xurl,hyperref}
\setlength{\parindent}{0pt}
\setlength{\parskip}{5pt}
\setlength{\emergencystretch}{3em}
\newcommand{\N}{\mathbb{N}}
\newcommand{\Z}{\mathbb{Z}}
\newcommand{\Q}{\mathbb{Q}}
\newcommand{\R}{\mathbb{R}}
\newcommand{\C}{\mathbb{C}}
\newcommand{\F}{\mathbb{F}}
\newcommand{\A}{\mathbb{A}}
\newcommand{\PP}{\mathbb{P}}
\newcommand{\E}{\mathbb{E}}
\newcommand{\eps}{\varepsilon}
\newcommand{\dd}{\,\mathrm{d}}
\begin{document}
\section*{UnsolvedMath ID: 5000002}
\textbf{Source problem code:} catalogue entry
\subsection{Definitions and mathematical statement}
Fix an integer \(n\ge5\) and a regular planar \(n\)-gon \(Q\) with distinguished vertex \(O\). Unfold a billiard trajectory by reflecting the current copy of \(Q\) across each side encountered, so the trajectory becomes a straight segment in the unfolded plane. The source calls a trajectory short when it joins two vertices and encounters no vertex in between, and calls a resulting endpoint vertex a reachable point.

For an oriented short trajectory the source uses endpoint angles \(\alpha,\beta\) from its Definition 2.1 to assign one of the \(n-2\) types
\[
 A_0,A_1,\ldots,A_{n-3},
\]
with indices understood modulo \(n-2\). The supplied information specifies only
\[
 \text{type }A_0\quad\Longleftrightarrow\quad\beta-\alpha=\frac{2\pi}{n}.
\]
The angle conventions and formulas for the remaining types are not included, so their meanings must remain those of the missing source definition.

A reachable polygon is an \(SL(2,\mathbb R)\) image of \(Q\) whose vertices other than \(O\) are reachable points. Here \(SL(2,\mathbb R)\) consists of real \(2\times2\) matrices of determinant \(1\); interpreting their linear action with \(O\) fixed uses coordinates with \(O\) at the origin, a normalization to be confirmed from the source. Let the vertices of any such polygon in clockwise order be
\[
 v_0=O,v_1,\ldots,v_{n-1}.
\]
Conjecture 2.3 asserts that their source-defined types satisfy
\[
 \operatorname{type}(v_j)=A_{j-1\bmod(n-2)}\qquad(1\le j\le n-1),
\]
that is, the ordered list is \(A_0,A_1,\ldots,A_{n-3},A_0\). This is a statement about every reachable polygon for every \(n\ge5\), not the existence of one polygon with the pattern.

A fully intrinsic formulation needs Definition 2.1: how the endpoint angles are measured, how the other \(A_k\) are assigned, the reference source vertex for reachability, and why an endpoint's type is well defined if there are multiple unfolding descriptions. None is guessed here. The supplied summary describes this as unresolved within the family of Conjectures 2.3--2.7 and reports no later resolution found through 2026; this is a supplied search assessment, not an independently verified current status.
\subsection{Short English statement}
Every reachable \(n\)-gon should have the specified cyclic sequence of vertex types, beginning and ending with \(A_0\), for all \(n\ge5\). The exact endpoint-angle and type definitions are missing from the excerpt and must be recovered from the cited source before the conjecture is fully formalized.
\subsection{Sources}
\begin{itemize}
\item Local registry: \url{data/unsolvedmath/problems.json}, permanent UnsolvedMath ID \(5000002\), source code catalogue entry; supplied statement, background, and research summary, updated 2026-08-06. Proposed by Dmitry Fuchs in 2020.
\item Dmitry Fuchs, \emph{Billiard Trajectories in Regular Polygons and Foliations by Closed Geodesics on Surfaces} (2020), Conjecture 2.3 and referenced Definition 2.1: \url{https://amj.math.stonybrook.edu/pdf-Springer-final/020-0170.pdf}. Supplied PDF-text extraction accessed 2026-07-29; no later authoritative resolution check is recorded.
\item The research summary mentions Fuchs (2017/2020) and Fuchs--Tabachnikov work on short trajectories and reachable points, but gives no additional bibliographic details.
\end{itemize}
\end{document}
